% R. P. Stanley, Some combinatorial properties of Jack symmetric functions,
% Advances in Mathematics 77 (1989), 76--115.
% DOI: 10.1016/0001-8708(89)90015-7
% Author's scan: https://math.mit.edu/~rstan/pubs/pubfiles/73.pdf
%
% Normalization (mathematical definitions of (2)--(4), Theorem 1.1,
% and Theorem 5.8 in self-contained notation).
\[
 z_\rho=\prod_{i\ge1}i^{m_i(\rho)}m_i(\rho)!,
 \qquad
 \langle p_\rho,p_\sigma\rangle
   =\delta_{\rho\sigma} z_\rho\alpha^{\ell(\rho)} .
\]
Here \(m_i(\rho)\) counts parts of size \(i\), \(\ell(\rho)\) counts
all parts, and \(p_\rho\) is the product of the corresponding power sums.
The \(J_\rho\) are mutually orthogonal, with
\[
 J_\rho=\sum_{\eta\le\rho}v_{\rho\eta}(\alpha)m_\eta,\qquad
 v_{\rho,1^n}=n! \quad(n=|\rho|).
\]
Here \(m_\eta\) is a monomial symmetric function, and \(\eta\le\rho\)
means \(|\eta|=|\rho|\) and
\(\sum_{i=1}^k\eta_i\le\sum_{i=1}^k\rho_i\) for every \(k\).
For a box \(s\) let \(A(s),l(s)\) be the number of boxes strictly
to its right and strictly below it, respectively, in English notation.
Write
\[
 H_\rho=\prod_{s\in\rho}(\alpha A(s)+l(s)+1),\qquad
 H'_\rho=\prod_{s\in\rho}(\alpha(A(s)+1)+l(s)),\qquad
 j_\rho=H_\rho H'_\rho .
\]
Thus \(J_\rho=H_\rho P_\rho\) for the monic Jack polynomial \(P_\rho\)
and \(j_\rho=\langle J_\rho,J_\rho\rangle\).

% Verbatim, Stanley, p. 110 (typesetting transcribed).
\paragraph{8.3. Conjecture.}
For fixed \(\lambda,\mu,\) and \(\nu\), the quantity
\(\langle J_\mu J_\nu,J_\lambda\rangle\) is a polynomial in
\(\alpha\) with nonnegative integer coefficients.

% Normalization cross-check: https://www.symmetricfunctions.com/jack.htm
% The scalar product in the conjecture is j_lambda times the coefficient
% of J_lambda in J_mu J_nu, NOT the coefficient itself.

% Additional source: Ryan Mickler, The Stanley conjecture revisited,
% arXiv:2401.01582v3, Sec. 3.1.1, Conjecture 3.3 (July 20, 2025).
% https://arxiv.org/pdf/2401.01582v3
% The following gives the precise equation in symbol rather than color
% notation. The arrays are read in source order from top to bottom,
% within each line from left to right; first mu, then nu, then lambda.
% Each row of the table specifies the coefficient on every word in it.
\[
\begin{array}{r|l}
 7 & \mathtt{UUUUUUUULUUL}\\
-2 & \mathtt{UUUUUUUUUUUL}\\
 1 & \mathtt{ULUUUUUUUUUL},\mathtt{UUUULUUUUUUL},\mathtt{UUUUUUUUULUL}\\
-2 & \mathtt{LUUUUUUULUUL},\mathtt{UULUUUUULUUL},\mathtt{UUULUUUULUUL},\\
   & \mathtt{UUUUULLULUUU},\mathtt{UUUUUUULLUUL},\mathtt{UUUUUUUULULL}\\
-2 & \mathtt{ULUUUUUULUUL},\mathtt{UUUULUUULUUL},\mathtt{UUUUUUUULLUL}\\
 1 & \mathtt{LLUUUUUULUUL},\mathtt{ULLUUUUULUUL},\mathtt{UUULLUUULUUL},\\
   & \mathtt{UUUULLLULUUU},\mathtt{UUUUUUULLLUL},\mathtt{UUUUUUUULLLL}\\
 1 & \mathtt{LUUUULLULUUU},\mathtt{UULLUUUULUUL},\mathtt{LUUUUUUULULL},\\
   & \mathtt{UULUUUULLUUL},\mathtt{UUULUUUULULL},\mathtt{UUUUULLLLUUU}.
\end{array}
\]
For \(\nu=(2,1)\) and three added boxes in distinct rows and columns,
the window \(R\) cuts out their three rows and three columns (taking
their intersections inside the partition). The arrays have shapes
\(21,21,321\) after deleting the unused rows and columns. Evaluate a word
\(D\) by taking the product of the indicated hooks of \(\mu,\nu\)
in its first six slots, divided by the product of the indicated hooks of
\(\lambda\) in its last six slots. Each hook is taken in the actual shape.
The factor \(F_R\) outside the window is the product of ratios of hooks
of \(\mu\) to hooks of \(\lambda\), using lower hooks in affected rows,
upper hooks in affected columns, and cancelling the unchanged hooks elsewhere.
Mickler asks that the sum solve the residue equations (46) for
all four spacings \({\bf v}\in{\mathbb Z}_{\ge0}^4\), and have weight zero.
The rational-function equality that implies those equations is
\[
 \frac{\langle J_\mu J_{21},J_\lambda\rangle}{j_\lambda F_R}
    =\sum_D c_D[D]_R .
\]
The weight of the sum is
\(\sum_D c_D(\#L_{\rm num}(D)-\#L_{\rm den}(D))\).
